\chapter{Revisão bibliográfica}
\label{ch:revisao}

Este capítulo apresenta os conceitos necessários para compreender o \CabotageLens{}. A revisão segue uma pergunta simples: o que precisa ser considerado para comparar, de forma justa, uma rota rodoviária direta com uma cadeia que utiliza cabotagem?

A resposta é construída em etapas. Primeiro, apresenta-se o contexto brasileiro. Depois, discutem-se mudança modal e comparação porta a porta. Em seguida, são explicadas as fontes ANTAQ e EU MRV, as fronteiras ambientais, as operações portuárias e a fronteira econômica. A literatura fornece o fundamento conceitual. Os resultados numéricos permanecem vinculados aos artefatos rastreados do projeto.

\section{Cabotagem brasileira, matriz de transporte e BR do Mar}
\label{sec:cabotagem-br}

Cabotagem é o transporte marítimo entre portos de um mesmo país. No Brasil, ela é relevante porque o território possui longas distâncias internas, extensa costa e forte dependência do transporte rodoviário. A agenda associada ao BR do Mar faz parte da discussão sobre diversificação da matriz de transporte e descarbonização dos corredores de carga \cite{icct2022}.

Esse contexto explica por que a cabotagem merece ser estudada. Ele não determina, por si só, que ela será melhor em toda rota. Por exemplo, uma origem distante de qualquer porto pode exigir um acesso rodoviário longo. Esse acesso pode reduzir parte do benefício esperado da navegação.

Portanto, o problema deve ser analisado rota a rota. As entradas relevantes são a origem, o destino, a carga e os portos. O processamento precisa incluir as distâncias terrestres e marítimas. A saída só é interpretável quando a unidade funcional e as fronteiras ambiental e econômica permanecem iguais nas alternativas comparadas.

\section{Mudança modal e short sea shipping}
\label{sec:sss-modal-shift}

Mudança modal é a transferência de uma carga de um modo de transporte para outro. No contexto deste trabalho, significa substituir parte de um percurso rodoviário por navegação costeira, sem deixar de considerar os acessos terrestres.

Revisões sobre \emph{short sea shipping}, ou navegação marítima de curta distância, mostram que essa transferência depende de custo, tempo, confiabilidade, integração terrestre, qualidade de serviço, características portuárias e barreiras institucionais \cite{modalshiftreview2020}. Estudos de emissões também indicam que o resultado varia com o corredor, a utilização e o tipo do navio, a distância, os acessos terrestres e as premissas de cálculo \cite{shortsea2019}.

A consequência para o \CabotageLens{} é direta. A pergunta não é ``o navio emite menos que o caminhão?''. A pergunta é: ``para esta remessa, nesta origem e neste destino, a cadeia completa com cabotagem apresenta melhor resultado que a rota rodoviária direta?''. A resposta vale para o cenário calculado e não para todos os corredores brasileiros.

\section{Comparações porta a porta e modelos de rede multimodal}
\label{sec:porta-a-porta}

Uma comparação porta a porta acompanha a carga desde a origem real até o destino real. Na alternativa multimodal, isso inclui o caminho até o porto de origem, a passagem pelos portos, a navegação e o caminho do porto de destino até o ponto final.

Um exemplo conceitual mostra a diferença. Se a rota rodoviária for comparada apenas com a distância navegada, os acessos aos portos desaparecem do cálculo. A alternativa marítima parecerá menor porque uma parte necessária do serviço foi omitida. A comparação correta soma todas as pernas que entregam a mesma remessa.

A literatura sobre competitividade da cabotagem conteinerizada e redes multimodais acrescenta outros elementos: frequência, tempo de espera, disponibilidade de serviço, capacidade, custo de inventário e componentes comerciais \cite{competitiveness2024}. O \CabotageLens{} não é um modelo comercial de super-rede. Ele não escolhe operadores, não confirma slots e não calcula tarifas de mercado. Seu papel é construir cenários auditáveis com a mesma unidade funcional e a mesma fronteira de cálculo.

\section{Dados operacionais da ANTAQ e indicadores do EU MRV}
\label{sec:antaq-mrv-literatura}

A ANTAQ disponibiliza duas tabelas complementares para este estudo. A tabela de Carga registra embarques, desembarques, massa, quantidade e portos relacionados à carga. A tabela de Atracação registra onde e quando o navio atracou e informa seu número IMO \cite{antaq2025}.

O número IMO é um identificador permanente do navio. Ele permite ligar as escalas observadas na ANTAQ aos indicadores anuais publicados pelo EU MRV. O MRV reúne informações de consumo e emissões de navios que operam no espaço regulado europeu, incluindo indicadores normalizados por trabalho de transporte \cite{eumrv2025}.

O processamento combina essas fontes em três etapas. Primeiro, ordena as escalas de cada viagem. Depois, reconstrói a carga a bordo entre as escalas. Por fim, busca no MRV a intensidade correspondente ao IMO. O valor MRV não é uma medição da viagem brasileira. Ele representa o desempenho anual reportado daquele navio e funciona como parâmetro para a atividade reconstruída na ANTAQ.

Nem todo IMO observado na ANTAQ possui correspondência no MRV. Nesses casos, o método utiliza um \emph{fallback}, isto é, uma regra alternativa previamente definida. O fallback por tipo usa aparagem simétrica das caudas quando a amostra permite e usa a mediana quando ela é pequena. O fallback por classe usa a média aparada registrada no artefato e recorre à mediana quando essa estatística não está disponível.

Essa regra reduz a influência de valores extremos sem escolher um número por conveniência. A saída identifica se a intensidade veio do IMO, da classe ou do tipo. Assim, dois resultados numericamente iguais podem ter níveis de evidência diferentes.

\section{Emissões operacionais, \ttw{}, \wtw{}, LCA, CO$_2$ e \cooe{}}
\label{sec:fronteiras-emissoes-literatura}

Uma fronteira ambiental define quais etapas entram no inventário de emissões. A sigla \ttw{} representa as emissões operacionais do uso do combustível: \emph{tank-to-wheel} na rodovia e \emph{tank-to-wake} na navegação. A abordagem \wtw{} também inclui as etapas anteriores ao uso do combustível. Estudos de LCA podem ampliar ainda mais a análise, incluindo produção de combustíveis, veículos, navios, infraestrutura e outros processos do ciclo de vida \cite{decarb2024,maritimelca2024}.

Essas fronteiras respondem a perguntas diferentes. Um valor \ttw{} não deve ser somado ou comparado diretamente com um valor \wtw{} sem alinhamento. O mesmo cuidado vale para CO$_2$, \cooe{} e CO$_2$eq: é necessário verificar quais gases, fatores e potenciais de aquecimento foram incluídos.

No \CabotageLens{}, a entrada ambiental é o consumo de combustível por perna e os fatores implementados. O processamento calcula emissões operacionais. A saída principal é \ttw{} em \cooe{}. Valores \wtw{}, LCA ou CO$_2$-only aparecem apenas como contexto, contraste ou possibilidade de trabalho futuro.

\section{Operações portuárias e hoteling}
\label{sec:portos-hoteling-literatura}

A cadeia multimodal não termina quando o navio chega ao porto. A permanência em berço pode exigir energia para sistemas auxiliares. Além disso, carga, descarga e atividades de terminal podem gerar consumo, custo e emissões \cite{berth2009,shipops2022,berthairquality2010}.

O termo \emph{hoteling} descreve o período em que o navio permanece no porto com sistemas de bordo em funcionamento. Operação portuária, por sua vez, abrange os componentes associados à movimentação e ao terminal representados no modelo. Os dois conceitos precisam ser distinguidos para evitar dupla contagem.

A inclusão desses componentes depende da proveniência. Um dado observado em um porto oferece evidência diferente de uma média entre portos ou de um valor genérico da literatura. Por isso, o sistema registra a fonte e a unidade. Quando não existe informação compatível, o componente permanece indisponível; ele não é convertido automaticamente em zero.

\section{Fronteira econômica e custo modelado}
\label{sec:custo-modelado-literatura}

O custo de uma operação logística real inclui preço contratado, tarifas, margens, seguro, confiabilidade, tempo, risco, inventário, frequência e disponibilidade de serviço \cite{competitiveness2024,modalshiftreview2020}. O \CabotageLens{} não representa todos esses elementos.

A saída monetária é um proxy de custo operacional modelado. Em termos simples, o programa combina os componentes de consumo e preço implementados no cenário. Essa saída é útil para uma comparação técnica sob regras comuns. Ela não é uma cotação, uma tarifa praticada ou um contrato de transporte.

Por exemplo, um cenário pode apresentar menor custo modelado e ainda enfrentar baixa frequência de serviço ou um preço comercial diferente. Portanto, o resultado econômico indica uma possibilidade dentro da fronteira do modelo. Ele não comprova competitividade comercial completa.

\section{Síntese da lacuna metodológica}
\label{sec:lacuna-metodologica}

A literatura sustenta três decisões deste trabalho. Primeiro, a cabotagem é relevante no contexto brasileiro, mas precisa ser avaliada por corredor. Segundo, a comparação deve ser porta a porta, porque acessos terrestres, portos e condições de serviço influenciam o resultado. Terceiro, emissões, custos e fontes de dados devem manter fronteiras e proveniência explícitas.

O \CabotageLens{} ocupa o espaço entre afirmações nacionais agregadas, comparações modais simplificadas e modelos comerciais de rede mais amplos. Sua entrada é um cenário específico. Seu processamento constrói as alternativas e preserva cada etapa. Sua saída permite verificar rotas, portos, distâncias, operações portuárias, emissões operacionais \ttw{} de \cooe{} e custo operacional modelado.

\begin{table}[htbp]
\centering
\small
\caption{Síntese do papel das principais famílias de literatura no relatório.}
\label{tab:uso-literatura}
\begin{tabularx}{\textwidth}{L{0.24\textwidth}Y Y}
\toprule
Família de literatura & Contribuição para o argumento & Leitura adotada no TF \\
\midrule
Cabotagem brasileira e BR do Mar & Situa o problema em um contexto nacional de matriz rodoviária, geografia costeira e política de cabotagem. & Motiva a análise rota a rota, sem converter contexto agregado em resultado por corredor.\\
Short sea shipping e mudança modal & Mostra que desempenho ambiental e adoção modal dependem de corredor, utilização, serviço e premissas. & Sustenta comparação condicional, porta a porta e sem generalização de superioridade modal.\\
Super-rede e competitividade & Evidencia a importância de rede, frequência, terminais, tempo, capacidade e componentes comerciais. & Delimita o \CabotageLens{} como triagem de cenário, não como otimização comercial completa.\\
WTW, LCA e combustíveis & Explica por que fronteiras ambientais e gases reportados alteram a interpretação. & Mantém a linha de base em \ttw{} de \cooe{} e usa outras fronteiras como contraste ou trabalho futuro.\\
Portos e hoteling & Indica que permanência em porto, consumo auxiliar e operações de terminal integram a cadeia multimodal. & Justifica incluir componentes portuários apenas com proveniência e fronteira compatíveis.\\
\bottomrule
\end{tabularx}
\end{table}

A Tabela~\ref{tab:uso-literatura} resume a função das fontes. Elas explicam por que a comparação deve ser brasileira, porta a porta, orientada por rota e transparente quanto às fronteiras. Os resultados, entretanto, são produzidos e classificados pelos artefatos do \CabotageLens{}.
